\section{Research Questions}

\subsection*{Main Research Question}

\begin{quote}
\itshape
How do reliability, system stability, and execution performance change when a workflow
orchestration system is migrated from a single-\vm{} Docker executor to a Kubernetes Run
Launcher under increasing concurrent workload, and at what specific workload level does this
migration become net beneficial as measured by the formally defined crossover point?
\end{quote}

\subsection*{Supporting Questions}

\begin{description}[leftmargin=2cm, style=nextline]

\item[\sq{1} --- VM Failure Threshold and Degradation Profile]
At what concurrency level do Docker containers running on a single \vm{} degrade in
reliability and performance, and how do job success rate, execution time variance, and
resource utilisation degrade as concurrent load increases toward that threshold?

\item[\sq{2} --- Pod Isolation Effectiveness and Blast Radius Containment]
To what degree does Kubernetes pod isolation contain failures and reduce execution time
variance compared to the single-\vm{} deployment under equivalent concurrent workload
levels?

\item[\sq{3} --- Operational Overhead and the Crossover Point]
What measurable scheduling latency and execution overhead does Kubernetes introduce compared
to the \vm{} Docker executor at each concurrency level, and at what concurrency level does
this overhead become smaller than the reliability gain from pod isolation --- the formally
defined crossover point?

\item[\sq{4} --- Workload Scalability Ceiling]
What is the maximum concurrent workload both architectures can sustain while maintaining a
job success rate above 95\%, and how do mean execution time and execution time variance
behave as concurrency approaches that ceiling?

\end{description}

\subsection*{How the Questions Form One Argument}

The four supporting questions are not independent. They form a single connected logical
chain.

\sq{1} establishes the exact failure profile of the \vm{} architecture --- the baseline
evidence. \sq{2} tests whether Kubernetes actually fixes what \sq{1} measured as broken.
\sq{3} synthesises \sq{1} and \sq{2} into the crossover point --- the primary contribution.
\sq{4} characterises the scalability ceiling of both architectures, framing the practical
limits of each deployment model.
